% SOURCE_AUTHORITY: sga5_fr_workpass.tex, lines 13120--13140 (LF-normalized unit).
% SOURCE_SHA256: 791F4EFFC5E02832D5D77ED03518C8156D6F07E4C8238B03545DB93D883FBB28
\begin{prooftext}[Demonstração de 7.14]
Seja $\pi$ o grupo fundamental de $W$. Então $F$ corresponde a um complexo de $A$-módulos $M^\bullet$ com grupo de operadores $\pi$. O complexo de $A$-módulos subjacente a $M^\bullet$ é isomorfo a $F_{\bar\eta}$ e, por conseguinte, pertence a $\Ob(\Delta)$. Introduzamos o grupo quociente $\bar\pi$ de $\pi$ por seu único $p$-subgrupo de Sylow $P$ [4, II 33], isto é, o limite projetivo de grupos quocientes discretos de $\pi$ cuja ordem é prima com $p$, onde $p$ é o expoente característico do corpo residual de $R$ (se $p=1$, põe-se $\bar\pi=\pi$). Tem-se, pois, a sequência exata
\[
1\longrightarrow P\longrightarrow \pi\longrightarrow \bar\pi\longrightarrow0,
\]
e obtém-se
\[
\RG_W(F)\simeq \R\Gamma^\pi(M^\bullet)
\simeq \R\Gamma^{\bar\pi}(\R\Gamma^P(M^\bullet))
\]
(«isomorfismo de Hochschild--Serre»). Ora, $P$ opera sobre $M^\bullet$ através de um grupo quociente finito, que é, portanto, um $p$-grupo e, por conseguinte, de ordem prima com $n$. Tem-se, pois,
\[
\R\Gamma^P(M^\bullet)\simeq \Gamma^P(M^\bullet),
\]
e $\Gamma^P(M^\bullet)$ é um somando direto de $M^\bullet$, pelo que pertence a $\Ob(\Delta)$. Por outro lado,
\[
\bar\pi=\prod_{\ell\ne p}\mathbb Z_\ell
\]
([4], II 33). Ficamos, pois, reduzidos a demonstrar que, se $M^\bullet\in\Ob(\Delta)$, então $\R\Gamma^{\bar\pi}(M^\bullet)\in\Ob(\Delta)$ e
\[
\clK{\Delta}(\R\Gamma^{\bar\pi}(M^\bullet))=0.
\]
